% TOP_PROBLEM: {"id": 5300026, "title": "Arc in a Cremer Julia set", "category_id": 11, "category": {"id": 11, "name": "dynamical_systems", "display_name": "Dynamical Systems", "description": "Problems about long-term behavior of deterministic systems, Hamiltonian dynamics, and periodic orbits.", "slug": "dynamical-systems", "order_index": 11, "created_at": "2026-07-31T15:26:25.671Z"}, "status": "open", "problem_number": "AMR-052-0026", "set_id": 13}
% FIRST_POSTED: 2026-10-02
% ENTRY_KIND: statement_only
% UnsolvedMath ID 5300026; source catalogue code AMR-052-0026
% !TeX program = lualatex
% Definition and sources only; this file is not an LLM solution attempt.
\documentclass[11pt,a4paper]{article}
\usepackage[margin=25mm]{geometry}
\usepackage{fontspec}
\setmainfont{lmroman10-regular.otf}[BoldFont=lmroman10-bold.otf,
 ItalicFont=lmroman10-italic.otf,BoldItalicFont=lmroman10-bolditalic.otf]
\usepackage{amsmath,amssymb,mathtools}
\usepackage{unicode-math}
\setmathfont{latinmodern-math.otf}
\AtBeginDocument{\let\setminus\smallsetminus}
\usepackage{xurl,hyperref}
\hypersetup{colorlinks=true,urlcolor=blue}
\setcounter{secnumdepth}{0}
\setlength{\parindent}{0pt}
\setlength{\parskip}{5pt}
\setlength{\emergencystretch}{3em}
\begin{document}
\section*{Arc in a Cremer Julia set}\label{5300026}
{\small UnsolvedMath ID 5300026 · AMR-052-0026 · Dynamical Systems · AMR Open Problem Lists}
\subsection{Definitions and mathematical statement}
Let $\alpha\in\mathbb R\setminus\mathbb Q$, put
$\lambda=e^{2\pi i\alpha}$, and consider $P_\alpha(z)=z^2+\lambda z$.
Assume the fixed point at zero is Cremer: there is no holomorphic
coordinate $h$ with $h(0)=0$, $h'(0)\ne0$, and
$h(P_\alpha(z))=\lambda h(z)$ on a neighborhood of zero.
Define the filled Julia set and Julia set by
\[
 K_\alpha=\{z\in\mathbb C:(P_\alpha^n(z))_{n\ge0}
                     \text{ is bounded}\},\qquad
 J_\alpha=\partial K_\alpha.
\]
The points $0$ and $-\lambda$ both lie in $J_\alpha$, and
$P_\alpha(-\lambda)=0$.

The question is whether every such quadratic Cremer polynomial admits
an arc in $J_\alpha$ joining these two points. Precisely, must there
exist a continuous injective map
\[
 \gamma:[0,1]\to J_\alpha,\qquad
 \gamma(0)=0,\quad\gamma(1)=-\lambda?
\]
Since the parameter interval is compact, its image is then homeomorphic
to a closed interval. The arc is not required to be smooth, rectifiable,
invariant, or an external ray.

Connectedness of the Julia set alone does not imply this conclusion,
because connected compact subsets of the plane need not be arcwise
connected. The question singles out the Cremer fixed point and its
other immediate preimage; it is weaker than asking that every pair of
Julia points be joined by an arc. Results for particular rotation
numbers therefore do not settle the universal formulation.
\subsection{Short English statement}
Must a quadratic Cremer Julia set contain an arc from the Cremer fixed point to its other immediate preimage?
\subsection{Sources}
\begin{itemize}
\item[S1] ULAM.AI and the UnsolvedMath Contributors, supplied problems.json, record 5300026 (AMR-052-0026), AMR Open Problem Lists. Catalogue identity and source transcription; CC BY 4.0.
\url{https://huggingface.co/datasets/ulamai/UnsolvedMath}.
\item[S2] Stony Brook - Open Problems in Dynamical Systems, 1992, Milnor local P1. Original source indexed by AMR-052-0026.
\url{https://www.math.stonybrook.edu/open-problems-dynamical-systems}.
\item[S3] Milnor, Problems on local connectivity, Problem1, in Bielefeld–Lyubich (1992).
\url{https://www.math.stonybrook.edu/preprints/ims92-7.pdf}.
\end{itemize}
\end{document}
